% Harvard-Format Professional CV — Daniel Saromo-Mori
% Compile with: xelatex cv-saromo.tex
% Premium tier: 2-page CV allowed

\documentclass[11pt,letterpaper]{article}

\usepackage[margin=0.6in, top=0.5in, bottom=0.5in]{geometry}
\usepackage{fontspec}
\usepackage{fontawesome5}
\usepackage[usenames,dvipsnames]{xcolor}
\usepackage{enumitem}
\usepackage{titlesec}
\usepackage{hyperref}
\usepackage{tabularx}

% --- Fonts (macOS system fonts) ---
\setmainfont{Times New Roman}
\setsansfont{Helvetica Neue}

% --- Colors (subtle, Harvard-appropriate) ---
\definecolor{headrule}{HTML}{2c3e50}
\definecolor{linkblue}{HTML}{1a5276}

% --- Hyperlinks ---
\hypersetup{colorlinks=true, urlcolor=linkblue, linkcolor=linkblue, pdfborder={0 0 0}}

% --- Section Formatting (Harvard style: centered, bold, with rules) ---
\titleformat{\section}
  {\large\bfseries\centering}
  {}
  {0em}
  {\vspace{-0.4em}\hrulefill\quad}[\quad\hrulefill\vspace{-0.4em}]
\titlespacing{\section}{0pt}{0.6em}{0.3em}

% --- List Formatting (tight bullets) ---
\setlist[itemize]{nosep, left=0pt, label=\textbullet, itemsep=1pt, topsep=2pt}

% --- No page numbers ---
\pagestyle{empty}

% --- Custom Commands ---
\newcommand{\entry}[4]{%
  \noindent\textbf{#1} \hfill #2 \par\nopagebreak
  \textit{#3} \hfill \textit{#4} \par\nopagebreak
}

\begin{document}

% ============================================
% HEADER
% ============================================
\begin{center}
  {\LARGE\bfseries Daniel Saromo-Mori} \\[4pt]
  \faIcon{map-marker-alt}\enspace Prague, Czechia
  \quad\faIcon{envelope}\enspace \href{mailto:danielsaromo@gmail.com}{danielsaromo@gmail.com}
  \\[2pt]
  \faIcon{linkedin}\enspace \href{https://linkedin.com/in/danielsaromo}{linkedin.com/in/danielsaromo}
  \quad\faIcon{github}\enspace \href{https://github.com/DanielSaromo}{DanielSaromo}
  \quad\faIcon{globe}\enspace \href{https://danielsaromo.xyz}{danielsaromo.xyz}
  \quad\faIcon{graduation-cap}\enspace \href{https://scholar.google.com/citations?user=ym6OfSgAAAAJ}{Google Scholar}
\end{center}

\vspace{-0.4em}
\noindent\rule{\textwidth}{0.8pt}
\vspace{-0.8em}

% ============================================
% EDUCATION
% ============================================
\section{Education}

\entry{Czech Technical University in Prague}{Prague, Czechia}{M.Sc. Cybernetics and Robotics, AI Specialization --- Currently Enrolled}{Sep 2024 -- Present}
\begin{itemize}
  \item Awarded merit scholarship for excellent study results in the Winter Semester of Academic Year 2024/2025
  \item Ranked 1st in Czechia in QS Engineering Ranking 2025
\end{itemize}
\vspace{0.3em}

\entry{Politecnico di Milano}{Milan, Italy}{M.Sc. Automation and Control Engineering --- Weighted Mark: 25.6/30}{Sep 2022 -- Feb 2024}
\begin{itemize}
  \item Completed coursework at institution ranked 13th in the world (QS Engineering 2022)
  \item Secured 1st place at Pitch Competition 2023 organized by Entrepreneurship Club POLIMI
\end{itemize}
\vspace{0.3em}

\entry{Pontificia Universidad Cat\'{o}lica del Per\'{u} (PUCP)}{Lima, Peru}{B.Sc. Mechatronics Engineering \& Professional Degree}{Mar 2014 -- Nov 2020}
\begin{itemize}
  \item Graduated in the top 6.66\% of the Faculty of Science and Engineering (6th of 32 mechatronics graduates)
  \item Thesis: ``Intelligent spider robot for detecting anti-personnel metallic landmines in uneven terrain''
  \item Earned Best Bachelor's Thesis award and Innovation Recognition Award at IMECE 2019 (world finals)
  \item Received Extraordinary Support Funding for Undergraduate Research Thesis; thesis rated outstanding unanimously
\end{itemize}
\vspace{0.3em}

% ============================================
% PUBLICATIONS
% ============================================
\section{Scientific Publications --- 200+ Citations (Google Scholar)}

\noindent\textbf{[C4]} Valdenegro-Toro, M. and \textbf{Saromo, D.} ``A Deeper Look into Aleatoric and Epistemic Uncertainty Disentanglement,'' \textit{LXCV Workshop at CVPR 2022}. New Orleans, USA. \textbf{200 citations.} Selected for oral presentation.\par\vspace{3pt}

\noindent\textbf{[C3]} Bravo, L., \textbf{Saromo, D.}, and Villota, E. ``Smart Insole Sensor for vGRF Measurement,'' \textit{9th International Symposium on Sensor Science}. Warsaw, Poland. 2022.\par\vspace{3pt}

\noindent\textbf{[C2]} \textbf{Saromo, D.}, Bravo, L., and Villota, E. ``Smart Sensor Calibration with Auto-Rotating Perceptrons,'' \textit{LXAI Workshop at ICML 2020}. Vienna, Austria. Selected for oral presentation.\par\vspace{3pt}

\noindent\textbf{[C1]} \textbf{Saromo, D.}, Villota, E., and Villanueva, E. ``Auto-Rotating Perceptrons,'' \textit{LXAI Workshop at NeurIPS 2019}. Vancouver, Canada. Selected for oral presentation.\par\vspace{3pt}

\noindent\textbf{[T1]} \textbf{Saromo, D.} ``Intelligent spider robot for detecting anti-personnel metallic landmines in uneven terrain,'' \textit{Pontificia Universidad Cat\'{o}lica del Per\'{u}}. Lima, Peru. 2020.\par\vspace{3pt}

\noindent\textbf{[Under Review]} ``Auto-Rotating Neural Networks: An Alternative Approach for Preventing Vanishing Gradients,'' \textit{Transactions on Machine Learning Research (TMLR)}.\par
\vspace{0.3em}

% ============================================
% RESEARCH EXPERIENCE
% ============================================
\section{Professional \& Research Experience}

\entry{Siemens}{Prague, Czechia}{Application Engineer, IT/OT Integration (Working Student)}{Aug 2025 -- Present}
\begin{itemize}
  \item Deliver IT/OT integration solutions at Siemens while concurrently pursuing M.Sc. in Cybernetics and Robotics at CTU Prague
\end{itemize}
\vspace{0.3em}

\entry{German Research Center for Artificial Intelligence (DFKI)}{Bremen, Germany (Remote)}{Guest Researcher}{Aug 2020 -- Dec 2020}
\begin{itemize}
  \item Extended the Auto-Rotating Perceptron (ARP) concept into a new neural network family (ARNN), implementing dense, recurrent, LSTM, GRU, and convolutional auto-rotating layers
  \item Validated ARNN performance against equivalent baseline models using the center's GPU clusters, presenting results at the Online Asian Machine Learning School (OAMLS) at ACML 2021
  \item Collaborated with research advisor Dr. Matias Valdenegro-Toro on architecture design and experimental methodology
\end{itemize}
\vspace{0.3em}

\entry{PUCP Applied Robotics and Biomechanics Research Group (GIRAB)}{Lima, Peru}{Research Assistant}{Mar 2020 -- Jul 2020}
\begin{itemize}
  \item Applied Auto-Rotating Perceptrons to wearable force sensor calibration, achieving 15x improvement in neural network performance (lower loss) compared to standard architectures
  \item Co-authored peer-reviewed paper presented at ICML 2020 under the guidance of research advisor Dr. Elizabeth Villota
\end{itemize}
\vspace{0.3em}

% ============================================
% TEACHING EXPERIENCE
% ============================================
\section{Teaching Experience}

\entry{PUCP Continuing Education Department}{Lima, Peru}{Lecturer --- Specialization Diplomas}{Sep 2019 -- Present}
\begin{itemize}
  \item Instructed AI for Games across 12 consecutive semesters (2019-2 through 2025-1), training 300+ professionals in applied AI techniques
  \item Taught Data Analysis Methods for Time Series in the Diploma in Data Analytics program (2022-1)
\end{itemize}
\vspace{0.3em}

\entry{PUCP Center for Advanced Manufacturing Technologies (CETAM)}{Lima, Peru}{Lecturer}{Sep 2020 -- Sep 2023}
\begin{itemize}
  \item Delivered ML for Industry across 7 semesters and Python for Data Science across 5 semesters, equipping manufacturing professionals with applied AI and data science capabilities
\end{itemize}
\vspace{0.3em}

\entry{National Meteorological and Hydrological Services (SENAMHI)}{Lima, Peru}{Lecturer --- Peruvian Government Entity}{May 2022 -- Jun 2022}
\begin{itemize}
  \item Designed and delivered introductory AI and ML curriculum for government meteorological professionals
\end{itemize}
\vspace{0.3em}

\entry{PUCP Undergraduate School}{Lima, Peru}{Teaching Assistant --- Faculty of Science and Engineering}{Mar 2019 -- Dec 2019}
\begin{itemize}
  \item Supported instruction in AI (2019-1), Machine Learning (2019-2), and Computer Science Applications (2019-2)
\end{itemize}
\vspace{0.3em}

% ============================================
% HONORS & AWARDS
% ============================================
\section{Honors \& Awards}

\begin{itemize}
  \item \textbf{Merit Scholarship}, Czech Technical University in Prague --- excellent study results (2024)
  \item \textbf{1st Place}, Pitch Competition 2023, Entrepreneurship Club POLIMI, Milan, Italy (2023)
  \item \textbf{LXAI Travel Grant} (\$1,860), NeurIPS 2019 --- oral and poster presenter, Vancouver, Canada (2019)
  \item \textbf{Innovation Recognition Award} (\$250), Old Guard 63rd Annual Oral Competition, World Finals at IMECE, Utah, USA (2019)
  \item \textbf{ASME Travel Award} (\$1,500), representing PUCP and South America at ASME IMECE, Utah, USA (2019)
  \item \textbf{1st Place + Technical Award} (\$850), Old Guard Oral Presentation Competition, ASME E-FEST South America, Lima (2019)
\end{itemize}
\vspace{0.3em}

% ============================================
% TALKS & PRESENTATIONS
% ============================================
\section{Talks \& Presentations --- 6 Countries}

\begin{itemize}
  \item \textbf{Oct 2024}: Auto-Rotating Neural Networks pitch, Dny AI series, \textit{Prague, Czechia}
  \item \textbf{Jun 2022}: Uncertainty Disentanglement paper, LXCV Workshop at CVPR, \textit{New Orleans, USA}
  \item \textbf{Nov 2021}: ARNN poster, Online Asian Machine Learning School at ACML, \textit{Singapore}
  \item \textbf{Jul 2020}: ARP sensor calibration paper, LXAI Workshop at ICML, \textit{Vienna, Austria}
  \item \textbf{Dec 2019}: Auto-Rotating Perceptrons paper, LXAI Workshop at NeurIPS, \textit{Vancouver, Canada}
\end{itemize}
\vspace{0.3em}

% ============================================
% SKILLS
% ============================================
\section{Skills \& Additional}

\begin{tabularx}{\textwidth}{@{}p{1.3in} X@{}}
\textbf{AI/ML:} & PyTorch, TensorFlow, Keras, Scikit-Learn, OpenCV, OpenAI Gym; Deep Learning (MLP, CNN, ARNN/ARP, DRL), Computer Vision, Transfer Learning \\[2pt]
\textbf{Robotics:} & ROS-1, Jetson Nano, Embedded Systems (Arduino, Raspberry Pi); TIA Portal (PLC/HMI), LabView \\[2pt]
\textbf{Programming:} & Python, MicroPython, C, C++, MATLAB, VBA, \LaTeX{}; NumPy, Pandas, Matplotlib, Plotly, Git, Linux \\[2pt]
\textbf{CAD/CAE:} & SolidWorks, Inventor, AutoCAD, EagleCAD, Altium Designer \\[2pt]
\textbf{Languages:} & Spanish (native), English (proficient), Italian (advanced), French (elementary) \\
\end{tabularx}

\end{document}
